% Grants and supervision as biblatex records (data model: uwhealthcv.dbx).
% Part of the template, loaded by uwhealthcv.sty after aprcv.sty. The CV
% keeps the records in .bib files and prints them with:
%
%   \cvowner{Name}               whose CV this is: \me in a grant prints it
%   \printgrants                 the RESEARCH GRANTS AND CONTRACTS table,
%                                newest first (by start, then end year)
%   \printsupervision{PhD}       one numbered list of students under a
%                                degree label: completed by date, then
%                                those in progress, by family name
%
% In the APR version a grant whose tenure overlaps the review period is
% highlighted as research (the whole table row); a student who completed in the
% period, or is in progress, is highlighted as teaching.

\RequirePackage{siunitx}

\def\cv@owner{}
\newcommand\cvowner[1]{\def\cv@owner{#1}}

\DeclareFieldFormat{cvplain}{#1}
\DeclareFieldFormat{cvcosupervisor}{\mkbibparens{co-supervised with #1}}

%% ----------------------------------------------------------------- grants
% The grants print as a longtable, like a hand-made table: headings repeat
% on each page, and a grant never splits across pages. biblatex cannot print
% entries into table cells, so the grant driver collects one table row per
% grant into \cv@rows (printing nothing), and \printgrants then sets them.
%
% Column widths: names, agency and short title wrap; amount and tenure don't.
\newlength\cvgrantnamewidth   \setlength\cvgrantnamewidth{4.5cm}
\newlength\cvgrantagencywidth \setlength\cvgrantagencywidth{4.6cm}
\newlength\cvgranttitlewidth  \setlength\cvgranttitlewidth{4.5cm}

% every investigator, never truncated to "et al." (biblatex's default
% truncates lists longer than three)
\ExecuteBibliographyOptions[grant]{maxitems=999,minitems=999}

% collect the investigators, one per line, into \cv@inv
\DeclareListFormat{cvcollect}{%
  \ifnumgreater{\value{listcount}}{1}{\gappto\cv@inv{\newline}}{}%
  \gappto\cv@inv{#1}}

% \cv@field{name}: a field's text, unexpanded (empty when it is not set)
\newcommand\cv@field[1]{%
  \ifcsundef{abx@field@#1}{}%
    {\unexpanded\expandafter\expandafter\expandafter{%
       \csname abx@field@#1\endcsname}}}

\gdef\cv@rows{}
\DeclareBibliographyDriver{grant}{%
  \gdef\cv@inv{}%
  \printlist[cvcollect]{investigators}%
  \aprsetentrycolour
  \ifdefempty\cv@rows{}{\gappto\cv@rows{\\}}% blank row between grants
  \xappto\cv@rows{%
    \noexpand\cv@grantline
      {\expandonce\aprentrycolour}%
      {\cv@field{role}}%
      {\expandonce\cv@inv}%
      {\cv@field{agency}}%
      {\cv@field{amount}}%
      {\cv@field{year}}%
      {\cv@field{endyear}}%
      {\cv@field{title}}}}

% \cv@grantline{colour}{role}{names}{agency}{amount}{year}{endyear}{title}
% \me is the CV's owner, in bold, with the role
\newcommand\cv@grantline[8]{%
  \ifx\relax#1\relax\else\rowcolor{#1}\fi
  \def\me{\textbf{\cv@owner\ifx\relax#2\relax\else\space(#2)\fi}}#3 &
  #4 &
  \$\num[group-separator={,},group-minimum-digits=4]{#5} &
  #6\ifx\relax#7\relax\else\ -- #7\fi &
  #8 \\*}

\DeclareSortingTemplate{cvgrants}{
  \sort[direction=descending]{\field{year}}
  \sort[direction=descending]{\field{endyear}\field{year}}
  \sort{\field{title}}
  \sort{\field{agency}}
}

\defbibenvironment{cvcollect}{}{}{}

\newcommand\printgrants{%
  \gdef\cv@rows{}%
  \setbox\z@\vbox{% collect the rows; nothing is printed here
    \begin{refcontext}[sorting=cvgrants]%
      \printbibliography[type=grant,env=cvcollect,heading=none]%
    \end{refcontext}}%
  \begin{longtable}[c]{>{\raggedright\arraybackslash}p{\cvgrantnamewidth}
                       >{\raggedright\arraybackslash}p{\cvgrantagencywidth}
                       l l
                       >{\raggedright\arraybackslash}p{\cvgranttitlewidth}}
  \underline{Names of Researcher(s)} & \underline{Agency} & \underline{\$ Amount} &
    \underline{Tenure} & \underline{Short Title} \\
  \endhead
  \cv@rows
  \end{longtable}}

%% ------------------------------------------------------------ supervision
\aprtypecategory{supervision}{teaching}

% full given names, whatever giveninits says for publications
\DeclareNameFormat{cvfull}{%
  \nameparts{#1}%
  \ifdefvoid\namepartgiven{}{\namepartgiven\space}%
  \ifdefvoid\namepartprefix{}{\namepartprefix\space}%
  \namepartfamily
  \ifdefvoid\namepartsuffix{}{\space\namepartsuffix}}

\newcommand\cv@monthname[1]{%
  \ifcase#1\or January\or February\or March\or April\or May\or June\or
  July\or August\or September\or October\or November\or December\fi}

% Name, date (or status), Title. Note (co-supervised with X) [Position]
\DeclareBibliographyDriver{supervision}{%
  \usebibmacro{begentry}%
  \printnames[cvfull]{author}%
  \setunit{\addcomma\space}%
  \iffieldint{year}
    {\printtext{%
       \iffieldint{month}{\cv@monthname{\thefield{month}}\space}{}%
       \thefield{year}}}
    {\printfield[cvplain]{status}}%
  \setunit{\addcomma\space}%
  \printfield[cvplain]{title}%
  \setunit{\addperiod\space}%
  \printfield[cvplain]{note}%
  \setunit{\space}%
  \printfield[cvcosupervisor]{cosupervisor}%
  \setunit{\space}%
  \printfield[brackets]{position}%
  \usebibmacro{finentry}}

\DeclareSortingTemplate{cvsupervision}{
  \sort{\field{year}\literal{9999}}
  \sort{\field[padside=left,padwidth=2,padchar=0]{month}\literal{99}}
  \sort{\name{author}}
}

% The list starts right after its degree label, so LaTeX treats the first
% \item specially and then clears \if@noparitem. biblatex typesets each
% entry in a group, which would undo that and treat every item as the
% first (the labels pile up), so clear it globally.
\newcommand\cv@degreeitem{\item\global\let\if@noparitem\iffalse}
\defbibenvironment{cvdegree}
  {\begin{enumerate}[label=\arabic*.]%
   \renewcommand*{\finentrypunct}{}}
  {\end{enumerate}}
  {\cv@degreeitem}

\def\cv@degree{}
\defbibcheck{cvdegree}{\iffieldequals{degree}{\cv@degree}{}{\skipentry}}

\newcommand\printsupervision[1]{%
  \begin{enumerate}
    \item[#1]%
    \def\cv@degree{#1}%
    \begin{refcontext}[sorting=cvsupervision]%
      \printbibliography[type=supervision,check=cvdegree,env=cvdegree,heading=none]%
    \end{refcontext}%
  \end{enumerate}}
